% ==============================================================================
% SECTION 5 — DISCUSSION
% ==============================================================================
\section{Discussion}
\label{sec:discussion}

\subsection{Hypothesis Evaluation}

\paragraph{H1 — Cooperative bias persists (\textcolor{green!60!black}{Confirmed}).}
The cooperative bias documented by Willis et al.\ is consistent across model generations
and providers. In the balanced noiseless condition---the most direct comparison
to the original benchmark---10 of 12 model--prompt combinations favour
cooperative equilibria as the plurality outcome ($\PC > \PA$ and $\PC > \PN$),
and aggressive equilibrium rates remain well below the 33\% prior for the
majority of configurations. The successor models (Claude~4.6, GPT-5.4~Mini)
show cooperative bias comparable to or stronger than their predecessors.
The within-lineage consistency is striking: Claude~4.6~Default (2\%A/50\%C)
and Claude~3.5~Sonnet~Default (4\%A/49\%C) are indistinguishable, suggesting
that the fundamental strategic disposition of Anthropic models is stable across
generations.

We hypothesise that this stability reflects convergence in alignment training
objectives: both models are trained with RLHF and Constitutional AI techniques
that strongly penalise defection-like behaviours. The persistence of cooperative
bias is consistent with this interpretation but does not rule out alternative
explanations (e.g., convergence in the training data distribution for
game-theoretic scenarios).

\paragraph{H2 — Aggressive capability parity (\textcolor{orange}{Partially Supported}).}
Evidence for H2 is mixed and prompt-dependent. The Self-Refine prompt increases
\ICD in all four models, replicating the finding from the original paper that
iterative self-critique improves aggressive strategy quality. Gemini~3.1~Pro~Refine
achieves the highest \ICD in the dataset (0.925), approaching payoff parity
between aggressive and cooperative agents; within the Anthropic lineage,
Claude~4.6~Refine (0.913) shows a larger narrowing than
Claude~3.5~Sonnet~Refine. This suggests that more capable models can
use self-refinement more effectively to generate strategic diversity.

However, under Default and Prose prompts, most models show \ICD values in the
range 0.45--0.82, not systematically higher than the original models. Moreover,
the models with the highest \ICD under Refine (Gemini~3.1~Pro at 0.925 and
Claude~4.6 at 0.913) are also those where the biased condition shows high
aggressive equilibria (62\% and 59\%A respectively),
suggesting that capability parity may genuinely increase the viability of
aggressive strategies. This finding reinforces the safety implication identified
by Willis et al.: prompt engineering techniques that improve aggressive
capabilities may inadvertently increase the risk of aggressive equilibria in deployed systems.

\paragraph{H3 — Cross-provider divergence (\textcolor{green!60!black}{Confirmed}).}
The four models exhibit substantially different evolutionary profiles.
Two-sample $z$-tests for proportions confirm that the cross-model differences
are not sampling artefacts. In the balanced noiseless condition (Default prompt),
3 of 6 pairwise comparisons of $\PA$ are significant at $p < 0.001$, in each case
contrasting Gemini~3.1~Pro (15\%) against the three models with near-zero aggressive
equilibria ($\leq 2\%$; $z \geq 3.30$). Under biased conditions, the effect sizes
are substantially larger: Gemini~2.5~Flash~Prose (70\%A) versus
Claude~4.6~Default (16\%A) yields $z = 7.71$ ($p < 0.0001$), and versus
GPT-5.4~Mini~Refine (24\%A) yields $z = 6.52$ ($p < 0.0001$).
Gemini~2.5~Flash is systematically the
most aggressive: its Prose and Refine prompts produce aggressive equilibrium
rates of 21--24\% in balanced conditions and 56--77\% in biased conditions.
GPT-5.4~Mini is the most cooperative: its Default and Refine prompts produce
aggressive equilibria of just 2--4\% in balanced conditions, and it achieves
the highest cooperative equilibrium rate in the dataset (70\% under Refine).
Claude~4.6 and Gemini~3.1~Pro occupy an intermediate position.

These cross-provider differences likely reflect distinct choices in training
data curation, RLHF reward modelling, and safety fine-tuning. The fact that
two Google models (Flash and Pro) differ substantially from each other---Flash
being more aggressive, Pro being more cooperative---suggests that model scale
and architecture matter independently of provider-level alignment decisions.

\paragraph{H4 — Noise robustness improves (\textcolor{red!70!black}{Suggestive, Not Robustly Confirmed}).}
The noise sensitivity of Claude~4.6 ($\Dnoise = 6$pp average) is lower than
Claude~3.5~Sonnet's ($\Dnoise = 13$pp), a 7pp difference in the direction of
improved robustness. However, a cross-study two-proportion test that propagates
the predecessor's smaller sample ($n = 100$) yields $z = 0.87$, $p = 0.385$:
the gap is directionally suggestive but not statistically robust, so we do not
claim confirmation.
Among the new models, Claude~4.6 is the most noise-robust (avg.\ 6pp),
followed by Gemini~3.1~Pro (8pp), while Gemini~2.5~Flash and GPT-5.4~Mini
are the most noise-sensitive (15pp each). GPT-5.4~Mini shows particularly
high sensitivity under the Refine prompt ($\Dnoise = 28$pp).

The inconclusive cross-study test for the Claude lineage is notable precisely because Willis et al.\
identified Claude's noise sensitivity as a key vulnerability---and although the
point estimate improves, substantial sensitivity persists under the Default and
Prose prompts ($\Dnoise = 9$ and $12$pp). One possible
explanation is that noise sensitivity is not a generalisation problem (which
better training would address) but rather a structural consequence of the
type of strategies generated: if cooperative strategies rely on clean mutual
cooperation signals (as Tit-For-Tat-like strategies do), noise will always
disrupt them, regardless of model capability.

\subsection{Implications for MAS Design}

The findings carry direct implications for the deployment of \LLM-based
multi-agent systems.

\paragraph{Provider choice matters.}
The choice of \LLM provider substantially
affects the emergent equilibrium dynamics of a deployed system. A system populated
with Gemini~2.5~Flash agents faces substantially higher risk of aggressive
equilibria than one using GPT-5.4~Mini or Claude~4.6.

\paragraph{Noise is a universal threat.}
Across all models and prompts,
action noise consistently shifts equilibria toward aggression in biased populations.
Systems deployed in noisy environments (e.g., with communication errors or
stochastic action execution) should be initialised with non-aggressive majorities,
or designed with noise-robust cooperation mechanisms.

\paragraph{Self-Refine as a double-edged tool.}
The Refine prompt generally
improves strategy quality but also narrows the gap between aggressive and
cooperative capabilities. Designers should be aware that using self-refinement
in strategy generation may inadvertently increase the viability of aggressive
equilibria.

\paragraph{Intra-generation stability is high.}
The near-identical results
between Claude~3.5~Sonnet and Claude~4.6~Default suggest that generational
model updates may not substantially alter cooperative biases when the alignment
training paradigm is unchanged.

\subsection{Limitations}

Estimation variance is the binding constraint on statistical precision: with
$n=100$ Moran iterations per condition, standard errors on equilibrium proportions
are non-trivial, and larger samples would sharpen inference. Second, the \IPD is a stylised social dilemma; results may not
generalise to more complex multi-player or asymmetric games. Third, we use
the same model for strategy generation and code conversion within each provider,
which may conflate generation capability with coding capability. Fourth, model
API versions may change over time, affecting reproducibility; we report exact
version identifiers in our codebase (available at \url{https://github.com/[AUTHOR]/evollm}).
% PRE-SUBMISSION: replace URL with actual repository before submission.
